\documentclass[11pt]{article}

\usepackage[margin=1in]{geometry}
\usepackage{microtype}
\usepackage{amsmath}
\usepackage{graphicx}
\usepackage{booktabs}
\usepackage[hidelinks]{hyperref}

\graphicspath{{figures/}}

\newcommand{\natex}{\texttt{natex}}

\title{An Honest Miss: Chinchilla as an Adoption Ramp in a Discontinuity Scan
of Tokens per Parameter}
\author{Hauke Hillebrandt\thanks{University College London. Email:
\texttt{ucjthhi@ucl.ac.uk}. This paper and the underlying \natex{} software were
prepared with substantial assistance from Anthropic's Claude models; the author
reviewed the analyses and text and is responsible for all remaining errors.}}
\date{July 2026 \\ \small analysis pass of record: 2026-07, \natex{} v0.2.0, seed 0}

\begin{document}

\maketitle

\begin{abstract}
\noindent A LoRD3 continuous-treatment discontinuity scan of $\log_{10}$
training tokens per parameter over calendar time, run blind on 549 language
models (2019--2024) from Epoch AI's notable-models panel, produces a single
significant top discovery: a break centered on 2022-09-22 (LLR $8.687$, scan
$p \leq 0.01$, robust to background degree and to neighborhood size
$k \geq 50$). The known truth --- the Chinchilla compute-optimal scaling
result, posted 2022-03-29, after which the global median tokens-per-parameter
ratio rose from $1.7$ to $29.2$ --- is \emph{not} localized: the best
candidate near the announcement ranks 158 of 547 (LLR $1.39$). The discovered
break is six months late and points the wrong way: a downward local contrast
of $1.63$ $\log_{10}$ units ($\times 43$) at the Sep/Oct 2022 boundary, where
the downstream edge of the Chinchilla adoption wave collides with a
measurement artifact --- fine-tuned releases whose recorded token counts
cover only the fine-tuning run over full base parameters. Nulls are reported
as nulls: the placebo test is vacuous (no non-forcing covariates), the
density test fails ($p=0.0041$) but is uninformative for a calendar-time
forcing, post-selection 2SLS is weak-instrument, and an independent kink
audit rejects a slope-change reading. The verdict is an artifact --- an
honest miss,
correctly diagnosed --- with two transportable lessons: diffusion-style
treatments need a ramp/slope statistic, not only a level-break statistic, and
fine-tune rows must be filtered from tokens-per-parameter panels before any
scan.
\end{abstract}

\section{Introduction}\label{sec:intro}

On 29 March 2022, Hoffmann et al.\ posted the ``Chinchilla'' result: for
compute-optimal training, model size and training tokens should scale in
equal proportion, implying that the then-standard practice --- descended from
the earlier scaling-law fits of \cite{kaplan2020} --- left frontier models
substantially undertrained on data \cite{hoffmann2022}. Release practice
responded within months, and by 2023 the industry had moved past
compute-optimality into deliberate over-training of small models for
inference economy, exemplified by LLaMA's trillion-token training runs
\cite{touvron2023} and rationalized by inference-adjusted scaling laws
\cite{sardana2024}.
Chinchilla itself has since been re-examined: \cite{besiroglu2024} replicate
its parametric fit and find inconsistencies in the reported estimates, while
the broader quantitative history of training-run inputs is documented by
\cite{sevilla2022} and maintained in Epoch AI's notable-models dataset
\cite{epoch2024}.

This gives a rare, clean testbed for \emph{automated natural-experiment
discovery}: a panel where the event, its date, its sign, and its approximate
magnitude are known ex ante. This paper runs the \natex{} toolkit's LoRD3
scan \cite{natex,herlands2018} on tokens per parameter over calendar time,
blind to the event date, and asks: does the scanner find Chinchilla? The
answer is no --- and the anatomy of the miss is the contribution.
Event-dated diffusion processes enter publication data as ramps, and a local
level-break statistic rewards sharp edges, not ramps; dating the cause is an
interrupted-time-series problem \cite{bernal2017}, not a discontinuity-scan
problem.

\section{Data}\label{sec:data}

The panel is a frozen extract of Epoch AI's notable-models database
\cite{epoch2024}: 549 language models released 2019-01 through 2024-12 with
non-missing parameter and training-token counts. Columns: release date (the
forcing variable, as days since 2019-01-01), $\log_{10}$ tokens per
parameter (the treatment), $\log_{10}$ parameters (the outcome), $\log_{10}$
training compute (21\% missing) and an open-weights flag (3.6\% missing),
both excluded from the scan to avoid listwise deletion. The known truth:
Chinchilla was posted 2022-03-29 (day $\approx 1183$); across that date the
global median tokens per parameter moves from $1.69$ to $29.17$ ($n = 185$
pre, $364$ post), a median-based shift of $+1.24$ $\log_{10}$ units
($+1.12$ on group means, $2.2 \to 29$; Mann--Whitney
$p = 1.6 \times 10^{-7}$ on an earlier extract, medians $2.2 \to 29.1$,
$n = 133/307$). At monthly resolution the shift is a ramp, not a step:
window medians run $1.7$ (pre-April 2022) $\to 5.7$ (mid-April to mid-May)
$\to$ ${\approx}28$ (mid-May to mid-July 2022).

\section{Design and Methods}\label{sec:methods}

\paragraph{LoRD3 continuous-treatment scan.} Following \cite{herlands2018} as
implemented in \natex{} v0.2.0 \cite{natex}, the scan fits a background model
of the treatment on the forcing variable (OLS, degree 1; degree 2 as
robustness), then scores every $k$-nearest-neighbor neighborhood ($k = 50$)
of every data point for the best split of its members into two half-spaces,
using a normal likelihood ratio on background residuals with locally
estimated variances. The top discovery is the neighborhood-split with the
maximum LLR; its $p$-value comes from a fitted-null Monte Carlo with $q = 99$
replicas and a $+1$-rank rule, so the minimum attainable $p$ is $0.010$ ---
reported throughout as $p \leq 0.01$, never as an exact value. The
specification is treatment $= \log_{10}$ tokens per parameter, forcing $=$
days since 2019-01-01 (forcing influence $1.0$), outcome $= \log_{10}$
parameters, no covariates; pipeline \texttt{natex study} $\to$
\texttt{natex discover --k 50 --q 99 --seed 0}, with localization detail and
robustness via the Python API, seed 0 throughout.

\paragraph{Validation battery, stated as run.} The battery's nulls and
refusals are reported, not patched over. (i)~Randomization: $p = 0.010$ at
both background degrees --- the $q=99$ floor. (ii)~Placebo covariates:
\texttt{placebo\_passed = True} but \emph{vacuous} ---
\texttt{placebo\_holm = \{\}}, because the design has no non-forcing
covariates to test. (iii)~Density: $p = 0.0041$, a \emph{fail}; this is
expected and uninformative here, because the forcing is calendar time with
strongly nonstationary publication intensity, and McCrary-style manipulation
logic \cite{mccrary2008} does not apply to a time forcing --- the test is
retained as a falsification check only. (iv)~Post-selection effects on
$\log_{10}$ parameters (2SLS and Wald at the discovered cutoff) are
\emph{advisory}: estimated on the discovery sample with no honest split and
instrumented by a partly artifactual break, they must not be read as causal
scaling-law parameters. (v)~An independent kink-design audit of the same
panel (a difference-in-slopes design \cite{bockerman2025}) cross-checks any
slope-change reading.

\section{Results}\label{sec:results}

\begin{figure}[t]
\centering
\includegraphics[width=\textwidth]{fig1}
\caption{(a) $\log_{10}$ tokens per parameter for all 549 models, with the
monthly median (line), the Chinchilla posting date (dashed), the scan's top
discovery (solid), and the seven verified fine-tune token-counting artifact
rows (open triangles; Table~\ref{tab:main}, panel~C). (b) LoRD3 scan profile:
each candidate neighborhood's LLR at its center date ($k=50$, degree 1,
$n=547$ candidates). The top discovery (filled marker, 2022-09-22, LLR
$8.69$, $p \leq 0.01$) sits at the Sep/Oct 2022 artifact edge; the best
candidate within $\pm 45$ days of the truth (open marker, 2022-04-15) scores
LLR $1.39$, rank 158/547. The 2023-06/07 cluster is the onset of the
deliberately over-trained era.}
\label{fig:fig1}
\end{figure}

\paragraph{The discovery.} The scan's top discovery (Table~\ref{tab:main},
panel A; Figure~\ref{fig:fig1}b) is centered on day 1360 $=$ 2022-09-22
(split normal $-1$ along the forcing axis): LLR $= 8.687$ ($8.687416$
unrounded in the run report), scan $p = 0.010$ --- the
$q=99$ Monte Carlo floor, i.e.\ $p \leq 0.01$ --- over $547$ candidate
neighborhoods from $n = 549$ rows. A degree-2 background reproduces the same
top center with LLR $8.836$ (randomization $p = 0.01$; null max LLRs
$1.8$--$7.6$ against observed $8.69$/$8.84$), and $k = 60$ moves it only to
2022-08-15; but $k = 30$/$40$ latch onto a 2020-10 composition boundary
instead (Table~\ref{tab:main}, panel B), so localization is stable only for
$k \geq 50$.

\paragraph{The known truth is not localized.} Within $\pm 45$ days of
2022-03-29 the best candidate neighborhood is centered 2022-04-15 with LLR
$1.39$ --- rank 158 of 547. The reason is visible in
Figure~\ref{fig:fig1}a: at daily resolution the Chinchilla shift enters the
published record as a 2--3 month adoption ramp, and a local level-break
statistic rewards sharp edges, not ramps.

\paragraph{What the discovered break actually is.} The top neighborhood
splits exactly at the Sep/Oct 2022 boundary into group A (2022-06-21 to
2022-09-22, $n=25$, median tokens/param $13.6$ --- the post-Chinchilla
adoption wave) versus group B (2022-10-03 to 2022-12-22, $n=25$, median
$0.4$): a \emph{downward} local group-mean contrast of $1.63$ $\log_{10}$
units ($\times 43$), against the known \emph{upward} global shift of $+1.12$
$\log_{10}$ ($+1.24$ median-based). Local-linear first-stage fits jump
$-1.7$ to $-2.0$ $\log_{10}$ at bandwidths of 120--240 days. Group B is
dominated by fine-tuned and continued-pretraining releases whose Epoch token
field records only fine-tuning tokens over full base parameters
(Table~\ref{tab:main}, panel C) plus small academic models. The top
discovery is therefore the adoption wave's downstream edge colliding with a
token-accounting artifact, ${\approx}6$ months after the event and with the
opposite sign. The secondary discovered cluster (2023-06-13 to 2023-07-06,
LLR $6.7$--$8.5$) is the onset of the deliberately over-trained era
\cite{touvron2023,sardana2024}; the December-2022 monthly median is already
$148$ tokens/param, and the median of the top-25 centers is day 1371 $=$
2022-10-03.

\paragraph{Effects and cross-checks, all null or advisory.} Post-selection
effects on $\log_{10}$ parameters at the discovered cutoff
(Table~\ref{tab:main}, panel D): the 2SLS estimate is weak-instrument
(first-stage $t = 1.89$; CI $[-0.747, 2.781]$); the Wald estimate is
$0.557$ (CI $[0.053, 1.062]$, first-stage $t = 4.01$). Both are
discovery-sample estimates with no honest split, instrumented by a partly
artifactual break --- diagnostics, not causal scaling-law parameters. The
independent kink audit \emph{rejects} a Chinchilla regime change read as a
slope change: a sign-unstable null with $t$ from $-0.91$ to $+1.48$ across
bandwidths.

\begin{table}[t]
\centering
\scriptsize
\caption{Chinchilla scan: numbers of record (\natex{} v0.2.0, seed 0).}
\label{tab:main}
\begin{tabular}{llll}
\toprule
Quantity & Value & Quantity & Value \\
\midrule
\multicolumn{4}{l}{\emph{Panel A: top discovery ($k=50$, $q=99$, degree 1; $n=549$, 547 candidates)}}\\
center & 2022-09-22 (day 1360) & LLR & $8.687$ \\
scan $p$ & $0.010$ = floor; $p \leq 0.01$ & local contrast & $-1.63\ \log_{10}$ ($\times 43$, down) \\
degree-2 center & 2022-09-22, LLR $8.836$ & first stage & $-1.7$ to $-2.0\ \log_{10}$ (bw 120--240d) \\
\midrule
\multicolumn{4}{l}{\emph{Panel B: $k$-sensitivity and truth localization}}\\
$k=30$ & 2020-09-29 (LLR $7.2$) & $k=40$ & 2020-10-21 (LLR $9.0$) \\
$k=50$ & 2022-09-22 (LLR $8.7$) & $k=60$ & 2022-08-15 (LLR $9.8$) \\
truth $\pm 45$d best & 2022-04-15, LLR $1.39$ & truth rank & 158 / 547 \\
global shift (true) & $+1.12\ \log_{10}$ ($2.2 \to 29$, up) & medians & $1.69 \to 29.17$ ($n=185/364$) \\
\midrule
\multicolumn{4}{l}{\emph{Panel C: verified fine-tune artifact rows (tokens/param; group B median $0.4$ vs.\ A $13.6$)}}\\
Flan-PaLM 540B & $0.0026$ & U-PaLM & $0.0024$ \\
LMSI-PaLM & $3.6 \times 10^{-6}$ & OPT-IML 175B & $0.011$ \\
Tk-Instruct & $0.095$ & BLOOMZ-176B & $0.114$ \\
mT0-13B & $1.54$ & & \\
\midrule
\multicolumn{4}{l}{\emph{Panel D: diagnostics and post-selection effects (advisory)}}\\
randomization & $p = 0.010$ (deg.\ 1 and 2) & placebo & passed, \emph{vacuous} (no covariates) \\
density & $p = 0.0041$, fails$^{\dagger}$ & kink audit & rejected; $t \in [-0.91, +1.48]$ \\
2SLS $\hat\tau$ & $1.017$ (se $0.900$), weak & CI & $[-0.747, 2.781]$, 1st-st.\ $t=1.89$ \\
Wald $\hat\tau$ & $0.557$ (se $0.257$) & CI & $[0.053, 1.062]$, 1st-st.\ $t=4.01$ \\
\bottomrule
\end{tabular}

\medskip
\parbox{0.95\textwidth}{\scriptsize Notes: $^{\dagger}$expected and
uninformative for a calendar-time forcing with nonstationary publication
intensity (McCrary manipulation logic does not apply). Scan $p$-values from a
$q=99$ fitted-null Monte Carlo with a $+1$-rank rule; $0.010$ is the minimum
attainable. Effects are estimated on the discovery sample with no honest
split. Artifact rows verified in the source panel: the token field records
fine-tuning tokens only, over full base parameters.}
\end{table}

\section{Caveats and Conclusion}\label{sec:conclusion}

Five caveats bound the claims. First, the value of this exercise is purely
methodological and diagnostic: no causal adoption claim is licensed; the scan
is a regime-edge detector, and dating the cause of a diffusion process needs
an interrupted-time-series or difference-in-differences framing
\cite{bernal2017} plus publication-lag domain knowledge. Second, the scan
$p = 0.010$ is the minimum attainable with $q = 99$ replicas and must be
reported as $p \leq 0.01$, not as an exact value. Third, the density test
fails ($p = 0.0041$) but is uninformative for a calendar-time forcing with
nonstationary publication intensity, and the placebo pass is vacuous
(\texttt{placebo\_holm = \{\}}; no non-forcing covariates to test). Fourth,
the post-selection 2SLS/Wald effects are estimated on the discovery sample
with no honest split and are instrumented by a partly artifactual break; they
must not be read as causal scaling-law parameters. Fifth, ``robust to $k$''
holds only for $k \geq 50$: at $k = 30$/$40$ the scan latches onto a 2020-10
composition boundary (tiny academic LMs against BERT-era corpus models), so
the robustness claim must be stated conditionally.

The verdict is an artifact --- an honest miss, correctly diagnosed: the scan
found real, significant structure whose two regimes bracket the truth, but it
dated the adoption wave's sharp downstream \emph{edge}, not the event, which
entered the record as a ramp. Two lessons transport beyond this dataset. For
method builders: scan batteries aimed at event-dated diffusion treatments
need a ramp/slope statistic --- a kink-style scan in the sense of
\cite{bockerman2025} --- alongside the level-break statistic, or
announcements will systematically lose to their own adoption edges. For
users of the Epoch panel \cite{epoch2024}: fine-tuned and
continued-pretraining rows manufacture a spurious low-tokens-per-parameter
cluster in late 2022 and must be filtered before tokens-per-parameter data
are used for scaling-law or event-study claims
\cite{hoffmann2022,besiroglu2024}.

\paragraph{Reproducibility.} All estimates: \natex{} v0.2.0 \cite{natex},
seed 0, on a frozen extract of the Epoch panel \cite{epoch2024} (not
committed). Figure~\ref{fig:fig1} regenerates deterministically from
\texttt{figures/make\_fig.py} (committed), which asserts the headline
numbers of record before drawing.

\scriptsize
\begin{thebibliography}{9}
\setlength{\itemsep}{2pt}\setlength{\parskip}{0pt}

\bibitem{bernal2017}
Lopez Bernal, J., Cummins, S., and Gasparrini, A. (2017).
Interrupted time series regression for the evaluation of public health
interventions: a tutorial.
\emph{International Journal of Epidemiology}, 46(1), 348--355.

\bibitem{besiroglu2024}
Besiroglu, T., Erdil, E., Barnett, M., and You, J. (2024).
Chinchilla scaling: A replication attempt.
arXiv:2404.10102.

\bibitem{bockerman2025}
B\"ockerman, P., Jysm\"a, S., and Kanninen, O. (2025).
\emph{Difference-in-Kinks Design}.
IZA Discussion Paper No.~18313.
\url{https://docs.iza.org/dp18313.pdf}

\bibitem{epoch2024}
Epoch AI (2024).
Data on notable AI models.
\url{https://epoch.ai/data/notable-ai-models}

\bibitem{herlands2018}
Herlands, W., McFowland~III, E., Wilson, A.~G., and Neill, D.~B. (2018).
Automated local regression discontinuity design discovery.
In \emph{Proceedings of the 24th ACM SIGKDD International Conference on
Knowledge Discovery and Data Mining (KDD~'18)}.

\bibitem{hoffmann2022}
Hoffmann, J., Borgeaud, S., Mensch, A., et al. (2022).
Training compute-optimal large language models.
In \emph{Advances in Neural Information Processing Systems}, 35 (NeurIPS 2022).
arXiv:2203.15556, posted 29 March 2022.

\bibitem{kaplan2020}
Kaplan, J., McCandlish, S., Henighan, T., et al. (2020).
Scaling laws for neural language models.
arXiv:2001.08361.

\bibitem{mccrary2008}
McCrary, J. (2008).
Manipulation of the running variable in the regression discontinuity design:
A density test.
\emph{Journal of Econometrics}, 142(2), 698--714.

\bibitem{natex}
Hillebrandt, H. (2026).
\emph{natex: automated natural-experiment discovery and estimation} (version
0.2.0). Software.
\url{https://github.com/HaukeHillebrandt/natex}

\bibitem{sardana2024}
Sardana, N., Portes, J., Doubov, S., and Frankle, J. (2024).
Beyond Chinchilla-optimal: Accounting for inference in language model scaling
laws.
In \emph{Proceedings of the 41st International Conference on Machine Learning
(ICML)}, PMLR 235, 43445--43460.

\bibitem{sevilla2022}
Sevilla, J., Heim, L., Ho, A., Besiroglu, T., Hobbhahn, M., and Villalobos, P.
(2022).
Compute trends across three eras of machine learning.
In \emph{2022 International Joint Conference on Neural Networks (IJCNN)}.

\bibitem{touvron2023}
Touvron, H., et al. (2023).
LLaMA: Open and efficient foundation language models. arXiv:2302.13971.

\end{thebibliography}

\end{document}
